\documentclass[10pt,a4paper]{amsart}
\usepackage[margin=19mm]{geometry}
\usepackage{amsmath,amssymb,mathtools}
\usepackage[scr=rsfs,cal=euler]{mathalfa}
\usepackage{xfrac}
\usepackage[unicode,hidelinks,bookmarks=false]{hyperref}
\newcommand{\ie}{i.e.\ }
\newcommand{\cf}{cf.\ }
\newcommand{\resp}{respectively\ }
\newcommand{\Subsection}{\subsection}
\providecommand{\nobreakdash}{\nobreak\hbox{-}}
\setlength{\emergencystretch}{2em}
\multlinegap0pt
\usepackage{amssymb}
\usepackage{bm}
\usepackage{mathtools}
\usepackage{algorithm}\usepackage{algpseudocode}
\usepackage{tikz}
\usepackage{xcolor}
\newtheorem{definition}{Definition}
\newtheorem{lemma}{Lemma}
\newcommand{\bn}{\boldsymbol{n}}
\newcommand{\btx}{\tilde{\bx}}
\newcommand{\bx}{\boldsymbol{x}}
\title{A Regularized Finite-Difference Approximation of Surface-restricted Emission and Reception Process in Acoustics with Application to Inverse Problems}
\author{Ashkan Javaherian}
\date{Source: arXiv:2212.04466v25 (CC BY 4.0)}
\begin{document}
\maketitle

\noindent\emph{Source and licence.} A Regularized Finite-Difference Approximation of Surface-restricted Emission and Reception Process in Acoustics with Application to Inverse Problems. Version arXiv:2212.04466v25 (CC BY 4.0), distributed by arXiv under the Creative Commons Attribution 4.0 International licence, http://creativecommons.org/licenses/by/4.0/, as recorded in the arXiv licence metadata for this exact version and shown on the abstract page https://arxiv.org/abs/2212.04466v25. Full licence notice: \texttt{licenses/arxiv-2212.04466-CC-BY-4.0.txt}.\par\smallskip

\noindent\textit{Editorial scope.} Counted unit: the article's lemma lemma:extraction (source0.tex lines 855--870). The article's complete original proof of it is delivered with it (source0.tex lines 873--929, one \texttt{proof} environment of 57 lines). No mathematics is written, repaired or re-derived: the statement and the proof are the article's own lines, in the article's own notation, and no retained line is substituted or re-flowed.  Retained uncounted as the source-local context the proof needs: lines 829--849.  Nothing was omitted from any retained span, so the package carries no omission record.  No retained line contains a citation, so the wrapper carries no bibliography and no bibliography entry.  No retained line contains a figure environment, an image-inclusion command or drawing code, so no image is delivered and the package declares no asset dependency.  Editorial formatting only, in addition: an AMS \texttt{amsart} preamble that restates the article's own declarations and macros used by the retained lines, namely the environments \texttt{definition}, \texttt{lemma} on separate counters, together with the standard packages those lines need.  The retained environments print in this document as Definition 1, Lemma 1 in order of appearance, whereas the article prints the same items with the numbering of its own full document.  The complete native source of the article as distributed in the arXiv e-print for this exact version is bundled unchanged as \texttt{sources/134648/source0.tex}.
\begin{definition}    \label{def:r}
We define a \emph{restriction operator} that, at each arbitrary but fixed time \( t \in [0, T] \), maps a field \( f \) in free space to the measured data \( y \) supported on \( \partial \Omega \). For each receiver \( r \), this operator is given by
\begin{align} \label{eq:r}
    \begin{split}
        &\mathcal{R}_r : C_0^{\infty} \big( \mathbb{R}^d \times [0,T] \big) \to C_0^\infty \big( \partial \Omega \times [0, T] \big), \\
        &\quad \mathcal{R}_r \left[ f \right](\bx_s, t) = y_r(\bx_s, t) :=
        \begin{cases}
            \displaystyle \frac{1}{2}
            \int_{\mathbb{R}^d} d\bx \ \delta_{b, \partial \nu_r}(\bx_s - \bx) \, f(\bx,t), & \text{if } (\bx_s, t) \in (\partial \nu_r \cap \partial \Omega) \times [0, T], \\
            0, & \text{otherwise}.
        \end{cases}
    \end{split}
\end{align}

\noindent
Here, \( \delta_{b,\partial \nu_r}(\bx_s - \bx) \) denotes a regularized Dirac delta distribution supported on the \((d{-}1)\)-dimensional infinite hyperplane \( \partial \nu_r  \). It is defined analogously to the standard Dirac delta distribution \( \delta_b \) in \( \mathbb{R}^d \), but with support restricted to the infinite plane (units: \( \mathrm{m}^{1-d} \)). It satisfies the sifting property
\begin{align} \label{eq:sifting-surface}
    f(\bx_s) = \lim_{b \to 0^+} \int_{\partial \nu_r } dS(\bx') \, \delta_{b,\partial \nu_r}(\bx_s - \bx') \, f(\bx'), \quad \bx_s \in \partial \nu_r \cap \partial \Omega.
\end{align}
The explicit form of the surface-restricted Dirac delta distribution is derived in the following Lemma.
\end{definition}

\begin{lemma} \label{lemma:extraction}
On the receiver surface, in the limit $b \to 0^+$, the field $f \in C_0^{\infty}\big(\mathbb{R}^d \times [0,T]\big)$ can be extracted by applying the restriction operator to its outward normal derivative in a narrow region surrounding the receiver surface. Specifically,
\begin{align}
    f(\bx_s, t)
    = \mathcal{R}_r \Big[ \frac{\partial f}{\partial \bn_{r,+}} (\bx, t) \Big],
    \qquad
    (\bx_s, t) \in \big(\partial \nu_r \cap \partial \Omega\big) \times [0,T],
\end{align}
where $\bn_{r,+}$ is the outward-pointing unit normal vector on $\partial \nu_r \cap \partial \Omega$, and $\mathcal{R}_r$ is defined in~\eqref{eq:r}. Moreover,
\[
    \delta_{b,\partial \nu_r}(\bx_s - \bx')
    = a_R(\bx_s) \, \delta_b(\bx_s - \bx'),
\]
where \( \delta_{b,\partial \nu_r} \) denotes the surface-restricted mollifier of bandwidth \( b \), and \( a_R(\bx_s) = \frac{d \bx_s}{dS(\bx_s)} \) is a scaling factor, which simplifies to \( a_R = b \) if one of the Cartesian coordinates coincides with the normal direction to the receiver surface.

\end{lemma}

\begin{proof}
At an arbitrary but fixed time \( t \in [0,T] \), the functions \( h, f \in C_0^{\infty} \big( \mathbb{R}^d \times [0,T] \big) \) satisfy:
\begin{align}  \label{eq:int-reception-n}
\begin{split}
     \int_{\partial \nu_r} dS \, h  f = \int_{ \nu_r^\pm} d\bx \, h \, \frac{\partial f}{\partial \bn_{r,\pm}} + \int_{\nu_r^\pm} d\bx \, \frac{\partial h}{\partial \bn_{r,\pm}}  f,
\end{split}
\end{align}
where $\bn_{r,+}$ and $\bn_{r,-}$ denote the outward and inward unit normal vectors, respectively, on the surface $\partial \nu_r \cap \partial \Omega$. We then choose the test function as
\begin{align}  \label{eq:delta-surface}
h(\bx) = \delta_b(\bx_s - \bx) +  \delta_b(\bx_s - \btx),
\end{align}
where \(\btx\) denotes the image point of \(\bx\) with respect to the infinite plane.

\noindent
Substituting Eq.~\eqref{eq:delta-surface} into both sides of 
Eq.~\eqref{eq:int-reception-n}, and summing the integral contributions from the two half-spaces, the second integral on the right-hand side vanishes because \(\tfrac{\partial h}{\partial \bn_{r,\pm}}=0 \). 


Additionally, in the specific setting that one Cartesian coordinate coincides with normal direction to \( \partial \nu_r \), and by 
decomposing the regularized Dirac delta distribution into sinc functions in the coordinates tangent and normal to the infinite plane in the left-hand side integral expression, and applying its 
sifting property in the limit \(b \to 0^+\) to localize the field and its 
normal derivative from the left-hand side integral expressions, we obtain
\begin{align} \label{eq:surface-integral2}
  \frac{1}{b} \,
  \Big[ f \big|_{\partial \nu_r^+} (\bx_s, t)- f \big|_{\partial \nu_r^-} (\bx_s, t)\Big]
  \;=\;\int_{\mathbb{R}^d} d\bx \ \delta_b(\bx_s - \bx) \, 
  \frac{\partial f}{\partial \bn_{r,+}}(\bx, t), \quad (\bx_s, t) \in \big( \partial \nu_r \cap \partial \Omega \big) \times [0,T],
\end{align}
which, in the limit \( b \to 0^+ \), indicates the emergence of a 
distributional singularity in 
\( \tfrac{\partial f}{\partial \bn_{r,+}} \) 
due to the jump of \( 2 f \) across the boundary, leading to the symmetric relation
\begin{align} \label{eq:surface-integral0}
  \lim_{b \to 0^+} \, \frac{2}{b}\, f\big|_{\partial \nu_r} (\bx_s, t)
  \;=\;\int_{\mathbb{R}^d} d\bx \ \delta_b(\bx_s - \bx) \, 
  \frac{\partial f}{\partial \bn_{r,+}}(\bx, t), \quad (\bx_s, t) \in \big( \partial \nu_r \cap \partial \Omega \big) \times [0,T].
\end{align}
For a general Cartesian coordinate system, performing the change of variables \( dS(\bx_s) \to d \bx_s \) in the integral on the left-hand side of Eq.~\eqref{eq:int-reception-n} yields the same formula as in Eq.~\eqref{eq:surface-integral0}, but with \( b \) replaced by the generalized scaling factor \( a_R(\bx_s) = \frac{d \bx_s}{dS(\bx_s)} \). This yields
\begin{align} \label{eq:surface-integral}
  \lim_{b \to 0^+} \, \frac{2}{a_R(\bx_s)}\, f\big|_{\partial \nu_r} (\bx_s, t)
  \;=\;\int_{\mathbb{R}^d} d\bx \ \delta_b(\bx_s - \bx) \, 
  \frac{\partial f}{\partial \bn_{r,+}}(\bx, t), \quad (\bx_s, t) \in \big( \partial \nu_r \cap \partial \Omega \big) \times [0,T].
\end{align}

\noindent
Eq.~\eqref{eq:surface-integral} yields the action of the restriction operator 
defined in Definition~\ref{def:r} on \(\tfrac{\partial f}{\partial \bn_{r,+}}\), 
where, in Eq.~\eqref{eq:r}, the surface-restricted Dirac delta distribution is 
given by
\begin{align}  \label{eq:dirac-delta-surface}
    \delta_{b,\partial \nu_r}(\bx_s - \bx) 
    = a_R(\bx_s)  \, \delta_b(\bx_s - \bx),\quad \bx_s \in \partial \nu_r \cap \partial \Omega ,
\end{align}
with \( a_R \) serving as a scaling factor representing the relative integral measure \( \frac{d \bx_s}{dS(\bx_s)} \). In the derivation above, we assumed the receiver surface to be an infinite plane. When a quantity is restricted to an infinite plane as a smeared jump—equivalently, as a smeared singularity in the normal derivative of the pressure—the restriction operator smooths the jump on the surface, i.e., it yields non-unique values that depend on the choice of \( a_R \). This non-uniqueness is related to the singularity that arises when inverting the double-layer potential integral formula in boundary element methods. From a numerical point of view, modeling the wave propagation and solving the forward problem uniquely yields the implicit formula~\eqref{eq:surface-integral}, rather than the explicitly surface-restricted quantity, since the latter depends on the factor \( a_R \).


\end{proof}

\end{document}
